\documentclass[a4paper, 12pt]{article}

\input{preamble.tex}

\lstset{
  language=Python,
  basicstyle=\ttfamily\small,
  numbers=left,
  numberstyle=\tiny,
  frame=single,
  breaklines=true,
  showstringspaces=false,
  captionpos=b,
  aboveskip=3mm,
  belowskip=3mm
}

\hypersetup{
  colorlinks=true,
  linkcolor=adblue,
  filecolor=adblue,
  urlcolor=adblue,
  citecolor=adblue,
}

\begin{document}

% =============================================
% ТИТУЛЬНЫЙ ЛИСТ
% =============================================
\begin{titlepage}
\centering
\vspace{0.75cm}
{\Huge\bfseries Разработка симулятора ADA4870\par}
\vspace{0.25cm}
{\Large для выходного канала трёхфазного высокочастотного генератора\par}
\vspace{0.25cm}
{\large (с предусилителями на ADA4807-2)\par}
\vspace{1.5cm}
{\large Техническая записка TN-2024-01\par}
\vspace{2cm}
\vfill
\begin{flushright}
Подготовил: grandfatherpikhto\\
Дата: \today
\end{flushright}
\vspace{1cm}
\end{titlepage}

\newpage
\tableofcontents
\newpage

% =============================================
% 1. ВВЕДЕНИЕ
% =============================================
\section{Введение и цель работы}

\subsection{Назначение схемы}

\par Это инженерный анализ схемы \texttt{ada4807\_4870.asc}~---
симуляции выходного каскада на основе операционного усилителя ADA4870
с предварительными усилителями тока ADA4807-2. Схема предназначена для
усиления сигнала токового выхода ЦАП AD9705 и питания низкоомной нагрузки

Основные характеристики усилителя ADA4870 \cite{ADA4870_ds}:
\begin{itemize}
  \item Частота единичного усиления: 52~МГц (большой сигнал),
        70~МГц (малый сигнал)
  \item Скорость нарастания выходного напряжения: 2500~В/мкс
  \item Максимальный выходной ток: 1~А
  \item Напряжение питания: от 10~В до 40~В (однополярное или
        двухполярное)
  \item Ток покоя: 32{,}5~мА (типовой)
\end{itemize}

Характеристики ADA4807-2 \cite{ADA4807_ds}:
\begin{itemize}
  \item Малошумящий rail-to-rail операционный усилитель
  \item Полоса пропускания 180~МГц (усиление +1)
  \item Скорость нарастания 225~В/мкс
  \item Низкий уровень шума: 3,1 нВ/√Гц
  \item Диапазон питания: 2,7~В … 10,5~В
\end{itemize}

\subsection{Состав файлов проекта}

\begin{itemize}
  \item \texttt{ada4807\_4870.asc} --- схема симуляции LTSpice
  \item \texttt{ada4807\_4870.net} --- список цепей (netlist)
  \item \texttt{ada4807\_4870.log} --- журнал Фурье-анализа
  \item \texttt{ada4807\_4870.raw} --- осциллограммы в формате CSV
\end{itemize}

\par Схема выглядит так:

\begin{figure}[H]
\centering
\begin{circuitikz}[american, scale=0.6, transform shape]
    %%%%%%%%%%%%%%%%%%%%%%%%%%
    % Верхний ОУ (Канал A)
    %%%%%%%%%%%%%%%%%%%%%%%%%
    \draw (0,14) node[op amp] (opampA) {}
        % 1. Питание и земля
        (opampA.+) -- ++(-0.5,0) node[ground] {} 
        (opampA.up) -- ++(0,0.3) node[above] {+2.5V}
        (opampA.down) -- ++(0,-0.3) node[below] {-2.5V}
        
        % 2. Инвертирующий вход -> Точка iouta -> Источник тока
        (opampA.-) to[short, -o] ++(-2,0) coordinate (iouta) node[left] {\(I_{outA}\)}
        (iouta) -- ++(0,-0.5) to[isource, l=$I_A$] ++(0,-1.5) node[ground]{}
        
        % 3. Возвращаемся к инвертирующему входу и рисуем ОС
        (opampA.-) -- ++(0,2) coordinate (fbA) 
        to[R, l={\(R_{fa}\)}, a=249R] (fbA -| opampA.out) -- (opampA.out)
        
        % 4. Рисуем выход прямо от ножки out
        (opampA.out) to[R, l={\(R_{b}\)}, a=100R, -o] ++(2.5,0) node[right] {\(OUT_A\)}; 
    
    % Общая подпись модели
    \node at ($(opampA) + (0, -2)$) {\textbf{ADA4807-2} (A)};  

    %%%%%%%%%%%%%%%%%%%%%%%%%%%
    % Нижний ОУ (Канал B)
    %%%%%%%%%%%%%%%%%%%%%%%%%%%
    \draw (0,7) node[op amp] (opampB) {}
        % 1. Питание и земля
        (opampB.+) -- ++(-0.5,0) node[ground] {} 
        (opampB.up) -- ++(0,0.3) node[above] {+2.5V}
        (opampB.down) -- ++(0,-0.3) node[below] {-2.5V}
        
        % 2. Инвертирующий вход -> Точка iouta -> Источник тока
        (opampB.-) to[short, -o] ++(-2,0) coordinate (ioutb) node[left] {\(I_{outB}\)}
        (ioutb) -- ++(0,-0.5) to[isource, l=$I_B$] ++(0,-1.5) node[ground]{}
        
        % 3. Возвращаемся к инвертирующему входу и рисуем ОС
        (opampB.-) -- ++(0,2) coordinate (fbB) 
        to[R, l={\(R_{fb}\)}, a=249R] (fbB -| opampB.out) -- (opampB.out)
        
        % 4. Рисуем выход прямо от ножки out
        (opampB.out) to[R, l={\(R_{b}\)}, a=88R7    , -o] ++(2.5,0) node[right] {\(OUT_B\)}; 

    \node at ($(opampB) + (0, -2)$) {\textbf{ADA4807-2} (B)};  


    %%%%%%%%%%%%%%%%%%%%%%
    % ОУ ADA4870-1
    %%%%%%%%%%%%%%%%%%%%%%
    \draw (10,10) node[op amp] (pwrAmp) {};
    \node at ($(pwrAmp) + (0, -2)$) {\textbf{ADA4870-1}};

    % Неинвертирующий вход (+) -> OUT_B
    \draw (pwrAmp.+) to[short, -o] ++(-1.0,0) node[left] {\(OUT_B\)};

    % Инвертирующий вход (-) -> OUT_A
    \draw (pwrAmp.-) to[short, -o] ++(-1.0,0) node[left] {\(OUT_A\)};

    % Цепь обратной связи (Rf || Cf)
    \draw (pwrAmp.-) -- ++(0,2) coordinate (fb_node) 
            node[circ] {}; % ТОЧКА СОЕДИНЕНИЯ У ВХОДА (Rf и линия от входа)
            
    \draw (fb_node) to[R, l={\(R_f\)}, a={1K4}] (fb_node -| pwrAmp.out) 
            coordinate (fb_end) node[circ] {} % ТОЧКА СОЕДИНЕНИЯ У ВЫХОДА (Rf и линия к выходу)
            -- (pwrAmp.out) node[circ] {}; % ТОЧКА НА ВЫХОДЕ ОУ

    % Конденсатор Cf (параллельно Rf)
    \draw (fb_node) -- ++(0,2) coordinate (cap_node) 
            % node[circ] {} % ТОЧКА ОТВЕТВЛЕНИЯ НА КОНДЕНСАТОР (слева)
            to[C, l={\(C_f\)}, a={1pF}] (cap_node -| fb_end) 
            % node[circ] {} % ТОЧКА ОТВЕТВЛЕНИЯ НА КОНДЕНСАТОР (справа)
            -- (fb_end);

    % Питание
    \draw (pwrAmp.up) to[short, -o] ++(0,0.5) node[above] {\(V_{DD}\)};
    \draw (pwrAmp.down)  to[short, -o] ++(0,-0.5) node[below] {\(V_{EE}\)};

    % Выходная часть
    \draw (pwrAmp.out) -- ++(1.5,0) coordinate (output_bus)
            node[circ] {}; % ТОЧКА СОЕДИНЕНИЯ С НАГРУЗКОЙ

    % Нагрузочный конденсатор об землю
    \draw (output_bus) to[C, l=$C_{load}$, a={\(220pF\)}] ++(0,-2) node[ground] {};

    % Выходная метка Signal
    \draw (output_bus) to[short, -o] ++(1.5,0) node[font=\boldmath, shift={(5pt,10pt)}] {Signal};

\end{circuitikz}
\caption{Общая схема подключения ADA4807-2 к ADA4870-1}
\end{figure}


\begin{itemize}
    \item Дифференциальную пару \(I_{outA}\), \(I_{outB}\) AD9705 выполняют два источника тока.
    \item Источники тока создают ток ~4mA (максимальный ток выхода AD9705 по документации не превышает 5mA)
    \item Паразитные параметры источников тока:
    \begin{itemize}
        \item Емоксть каналов \(I_{outA}\), \(I_{outB}\) --- 5pF
        \item Сопротивление каналов \(I_{outA}\), \(I_{outB}\) --- 150 \(k\Omega\)
    \end{itemize}
    \item Сигнал первого источника определяется формулой:
    \begin{align}
        I=\frac{I_{FS}}{2}  + \frac{I_{FS}}{2} \left( \sin{2\pi \cdot \omega \cdot t} \right)
    \end{align}
    \item Сигнал второго источника определяется формулой (сдвинут на \(\pi/2\)):
    \begin{align}
        I=\frac{I_{FS}}{2}  + \frac{I_{FS}}{2} \left( \sin{2\pi \cdot \omega \cdot t  + \pi} \right)
    \end{align}
    \item \(\omega\) --- частота. В данном случае, 1МГц
    \item \(V_{DD} = 18V\),
    \item \(V_{EE} = -18V\)
\end{itemize}

\par в LTSpice схема выглядит так:

\begin{figure}[H]
    \centering
    \includegraphics[width=0.9\linewidth, page=1]{scheme.png}
    \caption{Схема симуляции в LTSpice}
\end{figure}

% =============================================
% 2. ОПИСАНИЕ СХЕМЫ
% =============================================
\section{Описание схемы}

\subsection{Структурная схема}

Схема состоит из следующих функциональных блоков:
\begin{enumerate}
  \item Источники тока \(I_{OUTA}\) и \(I_{OUTB}\) --- модели токовых выходов
        ЦАП AD9705
  \item Входные трансимпедансные усилители (TIA) на сдвоенном ОУ ADA4807-2
        (U2, U3), преобразующие ток в напряжение
  \item Входная цепь ADA4870: резисторы \(R_{a}\), \(R_{b}\) и паразитные элементы
        \(R_{para}\)/\(C_{para}\) (модели выходного импеданса ЦАП)
  \item Дифференциальный усилитель ADA4870 (U1) с цепью обратной
        связи \(R_f\)/\(C_f\)
  \item Нагрузочная цепь: \(R_{Load}\) и \(C_{Load}\)
  \item Делитель напряжения \(R_1\)/\(R_2\) для управления входом SD
\end{enumerate}

\subsection{Источники тока \texorpdfstring{(\(I_{OUTA}\), \(I_{OUTB}\))}{{(IOUTA, IOUTB)}}}

Источники моделируют дифференциальный токовый выход ЦАП AD9705
\cite{AD9705_ds}. Токи изменяются в противофазе:
\begin{align}
  I_{\text{OUTA}} &= \frac{I_{FS}}{2}
    + \frac{I_{FS}}{2}\sin(2\pi f t),
    \label{eq:iouta}\\
  I_{\text{OUTB}} &= \frac{I_{FS}}{2}
    + \frac{I_{FS}}{2}\sin(2\pi f t + \pi),
    \label{eq:ioutb}
\end{align}
где $I_{FS}$ --- ток полной шкалы ЦАП (в текущей симуляции $I_{FS}=4$~мА,
что соответствует работе с пониженным выходным током),
$f = 1$~МГц --- частота сигнала.

Каждый источник является однонаправленным: ток изменяется
в диапазоне от 0 до $I_{FS}$, постоянная составляющая
$I_{FS}/2$ компенсируется дифференциальным включением.
В симуляции используются источники типа \texttt{BI}
(поведенческий источник тока LTSpice).

\subsection{Трансимпедансные усилители (ADA4807-2)}

Два канала сдвоенного ОУ ADA4807-2 (U2, U3) работают как
преобразователи ток--напряжение (TIA).
Схема включения каждого канала:
\begin{itemize}
    \item Неинвертирующий вход соединён с общим проводом (GND).
    \item Инвертирующий вход подключён к токовому выходу ЦАП
          (\(I_{OUTA}\) для U2, \(I_{OUTB}\) для U3).
    \item Резистор обратной связи \(R_3\) (249~Ом) включён между выходом
          и инвертирующим входом (для U2); \(R_5\) (249~Ом) для U3.
\end{itemize}
Виртуальная земля на инвертирующем входе обеспечивает
протекание входного тока через резистор обратной связи,
формируя выходное напряжение:
\begin{align}
    V_{N001} &= -I_{OUTA} \cdot R_3, \label{eq:v_n001}\\
    V_{N002} &= -I_{OUTB} \cdot R_5. \label{eq:v_n002}
\end{align}

Разностное напряжение на входах ADA4870 (ASA и ASB) равно
\begin{equation}
    V_{\text{diff}} = V_{N001} - V_{N002}
    = -(I_{OUTA} - I_{OUTB}) \cdot R_3
    = -I_{FS}\sin(2\pi f t) \cdot 249\,\Omega,
\end{equation}
поскольку \(I_{OUTA} - I_{OUTB} = I_{FS}\sin(2\pi f t)\).
Амплитуда дифференциального сигнала составляет
\[
    V_{\text{diff,amp}} = I_{FS} \cdot R_3 = 0{,}004\,\text{A} \times 249\,\Omega
    = 0{,}996\,\text{B}.
\]

Питание ADA4807-2 осуществляется от источников V3, V4: $\pm 2{,}5$~В.
Резисторы \(R_3\) и \(R_5\) выбраны равными 249~Ом из стандартного
ряда, обеспечивая высокую линейность преобразования при токе до 4~мА.

\subsection{Входная цепь ADA4870}

\subsubsection{Резисторы \texorpdfstring{\(R_{a}\) и \(R_{b}\)}{Ra и Rb}}

\textbf{\(R_{a}\)} (100~Ом) подключён между узлом N001 (выход U2, ASA)
и общим проводом (GND). Через него протекает выходной ток U2,
задавая потенциал узла ASA.

\textbf{\(R_{b}\)} (93{,}1~Ом) подключён между узлом N002 (выход U3, ASB)
и GND. Он обеспечивает симметрию импедансов для дифференциального
усилителя U1.

Условие симметрии (нулевая постоянная составляющая на выходе):
\begin{equation}
  \boxed{R_b = \frac{R_a \cdot R_f}{R_a + R_f}}
  \label{eq:sym}
\end{equation}

С новыми номиналами \(R_a=100\)~Ом, \(R_f=1400\)~Ом:
\[
    R_b = \frac{100 \times 1400}{100+1400} = \frac{140000}{1500}
    \approx 93{,}33\ \Omega.
\]
Выбран ближайший стандартный номинал 93{,}1~Ом.
Ошибка симметрии пренебрежимо мала, что подтверждается низкой
постоянной составляющей на выходе (порядка \(-0{,}5\)~В при
амплитуде сигнала 14~В).

\textbf{Важно:} перестановка \(R_{a}\) и \(R_{b}\) местами нарушила бы
условие, так как к узлу ASA дополнительно подключён резистор
обратной связи \(R_f\). Поэтому \(R_b\) всегда должно быть равно
\(R_a \parallel R_f\).

\subsubsection{Паразитные элементы ЦАП (\texorpdfstring{\(R_{paraA}\)/\(C_{paraA}\), \(R_{paraB}\)/\(C_{paraB}\)}{Rpara/Cpara})}

Эти элементы моделируют выходной импеданс токовых выходов ЦАП AD9705.
В симуляции приняты типовые значения из даташита \cite{AD9705_ds}:
\begin{itemize}
    \item \(R_{paraA} = R_{paraB} = 150\)~кОм,
    \item \(C_{paraA} = C_{paraB} = 5\)~пФ.
\end{itemize}
Параллельно им также включены конденсаторы Ca = Cb = 22~пФ (модель
паразитной ёмкости монтажа). Таким образом, суммарная ёмкость на каждом
входном узле (INN и INP) составляет \(5+22 = 27\)~пФ.

Влияние на низких частотах (1~МГц) незначительно. Полюс,
образованный ёмкостью 27~пФ и сопротивлением порядка 100~Ом,
лежит на частоте около 59~МГц, что выше рабочей полосы.

\subsection{Резисторы обратной связи \texorpdfstring{\(R_{fa}\), \(R_{fb}\)}{Rfa, Rfb}}
\par Значение 249 Ом для резисторов \(R_{fa}\) и \(R_{fb}\) в трансимпедансных усилителях (TIA) на ADA4807-2 выбрано не случайно, а исходя из баланса между несколькими ключевыми требованиями. Вот логика его расчёта:

\begin{itemize}
    \item Ток полной шкалы ЦАП AD9705 был установлен равным \(I_{FS} = 4 mA\) Этот ток и определяет амплитуду сигнала на выходе TIA.
    \item Целевое выходное напряжение: Для эффективного управления последующим каскадом на ADA4870, на входе которого для получения амплитуды ±14 В требуется дифференциальное напряжение около 1 В, было выбрано значение \(V_{diff_amp} \approx 0.996\)
    \item Для классического TIA на ОУ, где ток втекает в виртуальную землю инвертирующего входа, выходное напряжение определяется просто: \(V_{out} = I_{in} \times R_f\) В нашем случае:
\begin{align}
    R_f = \frac{V_{diff\_amp}}{I_{FS}} = \frac{0.996V}{0.004A} \approx 249 \Omega
\end{align}    
\end{itemize}

\par Таким образом, 249 Ом — это стандартный номинал из ряда E96, который идеально подошёл для преобразования тока 4 мА в напряжение, близкое к 1 В.
\par Можно было бы взять и другой номинал, но 249 Ом оказался оптимальным, поскольку:
\par Не перегружает ЦАП: Выходное напряжение ЦАП (около 1 В) остаётся в пределах его комплаенса (compliance voltage), что гарантирует линейность.
\begin{itemize}
    \item Не перегружает вход ADA4807: Выход TIA находится в пределах линейного диапазона ОУ при питании \(\pm 2.5V\).
    \item Оставляет запас по усилению: Даже небольшое изменение этого резистора привело бы к необходимости пересчёта \(Ra\) и \(R_b\)   в выходном каскаде для сохранения баланса.
\end{itemize}

\par Если нужно изменить выходной ток \(I_{FS}\) ЦАП, этот резистор нужно будет пересчитать по той же формуле, чтобы сохранить то же самое дифференциальное напряжение для каскада на ADA4870.

\subsection{Усилитель U1 (ADA4870)}

Усилитель ADA4870 включён по дифференциальной схеме.
Сигнал с выхода U2 (узел N001) подаётся на неинвертирующий вход
(ASA) через резистор \(R_a\), а сигнал с выхода U3 (узел N002) —
на инвертирующий вход (ASB) через \(R_b\). Обратная связь
организована резистором \(R_f\) с выхода AOUT на неинвертирующий
вход ASA.

Коэффициент усиления дифференциального сигнала:
\begin{equation}
    A_v = \frac{R_f}{R_a}.
    \label{eq:av}
\end{equation}

При \(R_f = 1{,}4\)~кОм и \(R_a = 100\)~Ом \(A_v = 14\).

Амплитуда выходного напряжения для заданного тока ЦАП:
\begin{equation}
    V_{\text{peak}} = A_v \cdot V_{\text{diff,amp}}
    = \frac{R_f}{R_a} \cdot I_{FS} \cdot R_3.
    \label{eq:vpeak}
\end{equation}

Для \(I_{FS}=4\)~мА, \(R_3=249\)~Ом, \(R_a=100\)~Ом, \(R_f=1{,}4\)~кОм
получаем:
\[
    V_{\text{peak}} = 14 \times 0{,}996 \approx 13{,}94\ \text{В}.
\]
Пиковый ток нагрузки \(R_{Load}=22\)~Ом:
\[
    I_{\text{peak}} = \frac{14\ \text{В}}{22\ \Omega} \approx 0{,}636\ \text{А} < 1\ \text{А}.
\]

Питание ADA4870: $\pm 18$~В от источников V1, V2.

\subsection{Цепь обратной связи \texorpdfstring{(\(R_f\), \(C_f\))}{Rf, Cf}}

\textbf{\(R_f\)} (1{,}4~кОм) — основной резистор обратной связи,
определяющий усиление.

\textbf{\(C_f\)} (2~пФ) включён параллельно \(R_f\). Образует нуль
в передаточной функции петли обратной связи, компенсирующий
полюс от паразитной ёмкости на входе ASA.
Частота нуля:
\[
    f_{z} = \frac{1}{2\pi R_f C_f}
    = \frac{1}{2\pi \times 1400 \times 2\times 10^{-12}}
    \approx 56{,}8\ \text{МГц}.
\]
Оптимальное значение \(C_f = C_{\text{INN}} \cdot R_a / R_f
\approx 27\ \text{пФ} \times 100 / 1400 \approx 1{,}9\ \text{пФ}\).
Используемое значение 2~пФ близко к расчётному и обеспечивает
устойчивость.

\subsection{Нагрузочная цепь \texorpdfstring{(\(R_{Load}\), \(C_{Load}\))}{Rload, Cload}}

\textbf{\(R_{Load}\)} = 22~Ом — активное сопротивление нагрузки.
Пиковый ток 0{,}636~А, что с запасом меньше максимального тока
ADA4870 (1~А).

\textbf{\(C_{Load}\)} = 220~пФ — модель паразитной ёмкости монтажа
и нагрузки. На частоте 1~МГц её импеданс около 723~Ом, что много
больше 22~Ом, поэтому шунтирующее действие минимально.

\subsection{Делитель SD (R1, R2)}

Делитель на резисторах R1 (2~кОм) и R2 (10~кОм) формирует
напряжение на входе разрешения SD усилителя ADA4870.
При питании $\pm 18$~В напряжение на SD составляет
\(-18\ \text{В} \times \frac{10}{12} = -15\ \text{В}\),
что соответствует диапазону включения (от \(-16{,}9\ \text{В}\)
до \(-13\ \text{В}\)).

% =============================================
% 3. ТЕОРЕТИЧЕСКИЙ РАСЧЁТ
% =============================================
\section{Теоретический расчёт}

\subsection{Алгоритм подбора номиналов}
\label{sec:algorithm}

На основе заданных требований (\(V_{out\_amp}\), \(R_{load}\), \(I_{FS}\)) и даташитов компонентов разработан автоматизированный алгоритм выбора \(R_a, R_f, R_b, C_f\). Ключевые этапы:

\begin{enumerate}
  \item \textbf{Проверка реализуемости.}
    \begin{itemize}
      \item Максимальная неискажённая амплитуда: \(V_{max} = |V_{sup}| - V_{headroom}\).
      \item Пиковый ток нагрузки: \(I_{peak} = V_{out\_amp}/R_{load}\). Если \(I_{peak} > I_{out\_max}\) или \(V_{out\_amp} > V_{max}\), расчёт прекращается.
    \end{itemize}
  \item \textbf{Дифференциальное напряжение TIA.}
    \(V_{diff\_amp} = I_{FS} \cdot R_{TIA}\) (обычно 0.996 В для 4 мА и 249 Ом).
  \item \textbf{Требуемое усиление.}
    \(A_v = V_{out\_amp} / V_{diff\_amp}\).
  \item \textbf{Перебор кандидатов \(R_a\).}
    Для каждого значения из списка \texttt{Ra\_candidates}:
    \begin{align*}
      R_f &= A_v \cdot R_a;\\
      R_b &= \frac{R_a \cdot R_f}{R_a + R_f}\quad \text{(идеальное условие баланса)};\\
      C_f &= C_{f\_base} + \frac{\max(0, R_f - R_{f\_threshold})}{1000} \cdot C_{f\_per\_kohm}.
    \end{align*}
    \(R_f\) и \(R_b\) округляются до ближайшего номинала из ряда E96 (1\%).
    Отбрасываются варианты с \(R_f > R_{f\_max}\).
  \item \textbf{Ранжирование вариантов.}
    Все допустимые комбинации сортируются по критерию \textbf{минимальной ошибки балансировки \(R_b\)}:
    \[
      \delta R_b = |R_{b\_ideal} - R_{b\_E96}| \to \min.
    \]
    Чем меньше \(\delta R_b\), тем меньше постоянная составляющая на выходе ADA4870. При равных \(\delta R_b\) предпочтение отдаётся комбинации, у которой \(R_f\) ближе к рекомендованному значению производителя (см. таблицу~\ref{tab:rf_rec}).
\end{enumerate}

\subsection{Рекомендованные значения \(R_f\) для ADA4870}
\label{sec:rf_recommendation}

ADA4870 — это усилитель с токовой обратной связью (CFB), и выбор \(R_f\) существенно влияет на полосу пропускания и устойчивость. Производитель указывает оптимальные \(R_f\) для типовых усилений (см. Table 6 даташита~\cite{ADA4870_ds}):

\begin{table}[H]
\centering
\caption{Рекомендованные \(R_f\) для ADA4870}
\label{tab:rf_rec}
\begin{tabular}{@{}cc@{}}
\toprule
\textbf{Усиление (V/V)} & \textbf{\(R_f\) (Ом)} \\
\midrule
+1  & 2000 \\
-1  & 1210 \\
+2  & 1500 \\
-2  & 1210 \\
+5, +10 & 1210 \\
\bottomrule
\end{tabular}
\end{table}

В нашей дифференциальной схеме эффективное усиление превышает 10, поэтому целевой ориентир для \(R_f\) составляет \textbf{1.21 кОм}. При сортировке комбинаций этот ориентир учитывается как вторичный ключ.

\subsection{Пример расчёта для данной схемы}

\begin{itemize}
    \item \(V_{\text{peak}} = 14\)~В, \(R_{Load} = 22\)~Ом $\Rightarrow I_{\text{peak}} = 0{,}636$~А.
    \item \(I_{FS} = 4\)~мА, \(R_{TIA} = 249\)~Ом $\Rightarrow V_{\text{diff,amp}} = 0{,}996$~В.
    \item \(R_a = 100\)~Ом, требуемое \(A_v = 14 / 0{,}996 \approx 14{,}06\).
    \item \(R_f = 14{,}06 \times 100 = 1406\)~Ом $\to$ выбираем 1{,}4~кОм (E96: 1,40) или 1,43~кОм (точнее).
        Примем \(R_f = 1400\)~Ом, тогда \(A_v = 14{,}00\).
    \item \(R_b = (100 \times 1400)/(1500) = 93{,}33\)~Ом $\to$ 93{,}1~Ом (E96).
    \item \(C_f \approx (27\times10^{-12}) \times 100 / 1400 \approx 1{,}9\)~пФ $\to$ оставляем 2~пФ.
    \item Проверка: \(V_{\text{peak}} = A_v \cdot V_{\text{diff,amp}} = 14 \times 0{,}996 = 13{,}94\)~В,
        что удовлетворяет требованию.
\end{itemize}

\subsection{Итоговые номиналы}

\begin{table}[H]
\centering
\caption{Номиналы компонентов схемы ada4807\_4870.asc}
\label{tab:components}
\begin{tabular}{@{}llp{6.5cm}@{}}
\toprule
\textbf{Компонент} & \textbf{Номинал} & \textbf{Назначение} \\
\midrule
\(R_f\)     & 1{,}4~кОм & Трансимпеданс, задаёт амплитуду \\
\(R_a\)     & 100~Ом   & Нагрузка на неинвертирующем входе U1 \\
\(R_b\)     & 93{,}1~Ом & Балансировочный резистор (\(R_a \parallel R_f\)) \\
\(R_3, R_5\) & 249~Ом   & TIA резисторы (преобразование ток--напряжение) \\
\(C_f\)     & 2~пФ     & Коррекция обратной связи \\
Ca, Cb     & 22~пФ    & Паразитная ёмкость входов ЦАП \\
\(R_{paraA}, R_{paraB}\) & 150~кОм & Выходное сопротивление ЦАП \\
\(C_{paraA}, C_{paraB}\) & 5~пФ   & Выходная ёмкость ЦАП \\
\(C_{Load}\)  & 220~пФ   & Паразитная ёмкость нагрузки \\
\(R_{Load}\)  & 22~Ом    & Нагрузка \\
R1, R2     & 2~кОм, 10~кОм & Делитель напряжения SD \\
\(V_{sup}\) & $\pm 18$~В & Питание ADA4870 \\
\(V_{sup,TIA}\) & $\pm 2{,}5$~В & Питание ADA4807-2 \\
\bottomrule
\end{tabular}
\end{table}

\subsection{Адаптивная настройка времени симуляции}
\label{sec:adaptive_tran}

Для получения достоверных результатов Фурье-анализа на любой частоте необходимо, чтобы в окне анализа укладывалось целое количество периодов сигнала и завершились переходные процессы. В программе реализована динамическая генерация директивы \texttt{.tran} на основе параметров, заданных в \texttt{config.json}:

\begin{itemize}
    \item \texttt{periods\_transient} = 10 – число периодов на затухание переходных процессов;
    \item \texttt{periods\_analysis} = 10 – число периодов, данные которых сохраняются для анализа;
    \item \texttt{points\_per\_period} = 1000 – минимальное количество точек на период.
\end{itemize}

По этим значениям автоматически вычисляются:
\begin{align}
    T &= \frac{1}{f_{\text{сигнала}}},\\
    t_{\text{end}} &= (N_{\text{transient}} + N_{\text{analysis}}) \cdot T,\\
    t_{\text{start}} &= N_{\text{transient}} \cdot T,\\
    \Delta t_{\text{max}} &= \frac{T}{N_{\text{points}}}.
\end{align}

Таким образом, для низких частот симуляция становится длиннее, но гарантирует, что переходные процессы успеют завершиться до начала сохранения, а Фурье-анализ выполняется на целом числе периодов. Это устраняет ложное поведение THD и обеспечивает физически достоверную зависимость искажений от частоты.

% =============================================
% 4. ТЕПЛОВОЙ РЕЖИМ
% =============================================
\section{Тепловой режим ADA4870}

Средняя мощность, рассеиваемая усилителем при синусоидальном
сигнале \cite{ADA4870_ds}:
\begin{equation}
  P_{\text{AVG}} = V_S \cdot I_q
    + \frac{2}{\pi}\cdot\frac{V_{CC}\,V_{\text{peak}}}{R_L}
    - \frac{V_{\text{peak}}^2}{2R_L},
  \label{eq:power}
\end{equation}
где \(V_S = 36\)~В (суммарное), \(I_q = 32{,}5\)~мА,
\(V_{CC} = 18\)~В, \(V_{\text{peak}} = 14\)~В, \(R_L = 22\)~Ом:
\[
    P_{\text{AVG}} = (36 \times 0{,}0325)
    + \frac{2}{\pi}\cdot\frac{18 \times 14}{22}
    - \frac{14^2}{2 \times 22}
    = 1{,}17 + \frac{2}{\pi}\cdot 11{,}4545 - 4{,}4545
    \approx 1{,}17 + 7{,}29 - 4{,}45 = 4{,}01\ \text{Вт}.
\]

Температура перехода при окружающей среде \(T_{\text{amb}} = 40\)~°C
и тепловом сопротивлении корпус–среда \(\theta_{CA}\):
\[
    T_J = T_{\text{amb}} + P_{\text{AVG}} \cdot (\theta_{JC} + \theta_{CA}),
\]
где \(\theta_{JC} = 1{,}3\)~°C/Вт (для корпуса SOIC).
Для гарантированного запаса (\(T_J \le 150\)~°C) требуется
\[
    \theta_{CA} \le \frac{150 - 40}{4{,}01} - 1{,}3 \approx 27{,}4 - 1{,}3 = 26{,}1\ \text{°C/Вт}.
\]
Рекомендуется радиатор с тепловым сопротивлением не более 20~°C/Вт.

Мощность на нагрузке:
\[
    P_{R_{\text{load}}} = \frac{V_{\text{peak}}^2}{2R_{\text{load}}}
    = \frac{196}{44} \approx 4{,}45\ \text{Вт}.
\]
Нагрузочный резистор должен быть рассчитан на мощность не менее 10~Вт.

% =============================================
% 5. РЕЗУЛЬТАТЫ МОДЕЛИРОВАНИЯ
% =============================================
\section{Результаты моделирования в LTSpice}

Моделирование выполнено в LTSpice версии 26.0.1 для Windows \cite{LTspice}.
В базовой симуляции использовалось значение \(R_f = 820\)~Ом, дающее
амплитуду около \(8{,}65\)~В на нагрузке 22~Ом.
Директивы симуляции:
\begin{itemize}
  \item \texttt{.tran 10us}
  \item \texttt{.FOUR 1Meg V(aout)}
  \item \texttt{.MEAS TRAN Vfinal FIND V(aout) AT=10us}
\end{itemize}

\subsection{Параметры симуляции}

\begin{table}[H]
\centering
\caption{Параметры (.param) симуляции}
\label{tab:params}
\begin{tabular}{@{}lll@{}}
\toprule
\textbf{Параметр} & \textbf{Значение} & \textbf{Назначение} \\
\midrule
IFS   & 4~мА   & Ток полной шкалы ЦАП \\
FREQ  & 1~МГц  & Частота сигнала \\
\(R_f\)Val & 820~Ом & Резистор обратной связи (базовая конфигурация) \\
\(R_{Load}\)Val & 22~Ом & Сопротивление нагрузки \\
\bottomrule
\end{tabular}
\end{table}

\subsection{Фурье-анализ выходного сигнала (Rf = 820 Ом)}

\begin{table}[H]
\centering
\caption{Фурье-анализ V(aout), \(f_0 = 1\)~МГц (из файла .log)}
\label{tab:fourier}
\begin{tabular}{@{}rrrrrr@{}}
\toprule
\textbf{№} &
\textbf{Частота, Гц} &
\textbf{Амплитуда, В} &
\textbf{Норм.} &
\textbf{Фаза, °} &
\textbf{Норм. фаза, °} \\
\midrule
1 & \(1{,}000\times10^6\) & \(8{,}645\times10^0\) & \(1{,}000\) & \(92{,}73\) & \(0{,}00\) \\
2 & \(2{,}000\times10^6\) & \(1{,}647\times10^{-4}\) & \(1{,}905\times10^{-5}\) & \(140{,}79\) & \(48{,}06\) \\
3 & \(3{,}000\times10^6\) & \(1{,}254\times10^{-4}\) & \(1{,}450\times10^{-5}\) & \(-69{,}62\) & \(-162{,}35\) \\
4 & \(4{,}000\times10^6\) & \(1{,}373\times10^{-4}\) & \(1{,}588\times10^{-5}\) & \(94{,}03\) & \(1{,}30\) \\
5 & \(5{,}000\times10^6\) & \(1{,}213\times10^{-4}\) & \(1{,}403\times10^{-5}\) & \(-127{,}85\) & \(-220{,}57\) \\
6 & \(6{,}000\times10^6\) & \(1{,}266\times10^{-4}\) & \(1{,}465\times10^{-5}\) & \(37{,}79\) & \(-54{,}94\) \\
7 & \(7{,}000\times10^6\) & \(1{,}439\times10^{-4}\) & \(1{,}665\times10^{-5}\) & \(179{,}05\) & \(86{,}32\) \\
8 & \(8{,}000\times10^6\) & \(1{,}378\times10^{-4}\) & \(1{,}594\times10^{-5}\) & \(-13{,}93\) & \(-106{,}65\) \\
9 & \(9{,}000\times10^6\) & \(1{,}650\times10^{-4}\) & \(1{,}908\times10^{-5}\) & \(137{,}08\) & \(44{,}35\) \\
\bottomrule
\end{tabular}
\end{table}

\begin{table}[H]
\centering
\caption{Сводка результатов симуляции (Rf = 820 Ом)}
\label{tab:summary}
\begin{tabular}{@{}lr@{}}
\toprule
\textbf{Параметр} & \textbf{Значение} \\
\midrule
Амплитуда 1-й гармоники       & 8{,}645~В \\
Постоянная составляющая (DC)  & \(-0{,}5446\)~В \\
Полный THD                    & 0{,}00633\% \\
\(V_{\text{aout}}\) при \(t=10\)~мкс & \(-0{,}956\)~В \\
\(V_{SD}\)                      & \(-15{,}183\)~В \\
\bottomrule
\end{tabular}
\end{table}

Примечание: измеренный DC-сдвиг (\(-0{,}54\) В) связан с тем, что
в данной симуляции резистор \(R_b\) был выбран 88{,}7~Ом
(под \(R_f=820\)~Ом). При изменении \(R_f\) на 1{,}4~кОм и \(R_b\)
на 93{,}1~Ом постоянная составляющая уменьшается.

\subsection{Ожидаемые характеристики при Rf = 1,4 кОм}

При увеличении \(R_f\) до 1{,}4~кОм и балансировке схемы
ожидаются следующие параметры:
\[
    V_{\text{peak}} \approx 14\ \text{В},\quad
    \text{DC offset} \approx -0{,}5\ \text{В},\quad
    \text{THD} < 0{,}01\%.
\]
Пиковый ток в нагрузке 22~Ом составит около 0{,}64~А.
Тепловыделение ADA4870 возрастёт до \(\sim 4\)~Вт,
что требует радиатора.

\section{Результаты симуляции}

\par На выходах \(I_{NN}\) и \(I_{NP}\) получаем:
\begin{figure}[H]
\centering
\begin{tikzpicture}
\begin{axis}[
    width=12cm, height=7cm,
    xlabel={Время (с)},
    ylabel={Напряжение \(V(inn)\) (В)},
    grid=major,
    legend pos=north east,
    % Настройка разделителя (для CSV это запятая)
    table/col sep=comma 
]
    % График для V(inn)
    \addplot[red, thick] table [
        x index=0, 
        y index=2,
        % each nth point=10 % Раскомментируйте, если файл очень тяжелый
    ] {./data/data.csv};
    \addlegendentry{$V(inn)$}

    % График для V(inp)
    \addplot[blue, thick] table [
        x index=0, 
        y index=3
    ] {./data/data.csv};
    \addlegendentry{$V(inp)$}
    
    \legend{\(V_{inn}\), \(V_{inp}\)}
\end{axis}
\end{tikzpicture}
\caption{График напряжения на узлах \(I_{NN}\), \(I_{NP}\)}
\end{figure}

\par Тогда, выходной сигнал будет выглядеть так:

\par На выходах \(I_{NN}\) и \(I_{NP}\) получаем:

\begin{figure}[H]
\centering
\begin{tikzpicture}
\begin{axis}[
    width=12cm, height=7cm,
    xlabel={Время (с)},
    ylabel={Напряжение \(V_{signal}\) (В)},
    grid=major,
    legend pos=north east,
    % Настройка разделителя (для CSV это запятая)
    table/col sep=comma 
]
    % График для V(signal)
    \addplot[green!50!black, thick] table [
        x index=0, 
        y index=13] {./data/data.csv};

    \addlegendentry{$V_{signal}$}
   
    \legend{\(V_{signal}\)}
\end{axis}
\end{tikzpicture}
\caption{График напряжения на выходе ADA4870-1 \(V_{signal}\)}
\end{figure}

\begin{figure}[H]
\centering
\begin{tikzpicture}
\begin{axis}[
    width=12cm, height=7cm,
    xlabel={Время (с)},
    ylabel={Ток \(R_{load}\) (mA)},
    grid=major,
    legend pos=north east,
    % Настройка разделителя (для CSV это запятая)
    table/col sep=comma 
]
    % График для I(load)
    \addplot[brown, thick] table [
        x index=0, 
        y index=31] {./data/data.csv};

    \addlegendentry{$V_{signal}$}
   
    \legend{\(I_{load}\)}
\end{axis}
\end{tikzpicture}
\caption{График тока на \(R_{load}\)}
\end{figure}

\subsection{Анализ деградации THD от частоты}
\label{sec:degradation}

Для оценки рабочего диапазона частот и поведения нелинейных искажений применяется методика \texttt{degradation\_sweep} – последовательный запуск симуляции на нескольких частотах с фиксированными номиналами схемы. Чтобы исключить артефакты, связанные с переходными процессами и недостаточной длительностью анализа, используется адаптивное задание времени симуляции (см. раздел~\ref{sec:adaptive_tran}).

На рисунке~\ref{fig:thd_vs_freq} приведён график зависимости THD от частоты, полученный в диапазоне 100\,кГц – 5\,МГц для выбранной комбинации номиналов (\(R_a=100\)\,Ом, \(R_f=1400\)\,Ом, \(R_b=93.1\)\,Ом). Виден монотонный рост искажений с частотой от \(\sim 0.048\%\) до \(\sim 0.058\%\), что характерно для усилителей с токовой обратной связью при уменьшении глубины обратной связи.

\begin{figure}[H]
\centering
\begin{tikzpicture}
\begin{axis}[
    width=12cm, height=7cm,
    xlabel={Частота (Гц)},
    ylabel={THD (\%)},
    grid=major,
    legend pos=north west,
    xmode=log,
    log basis x={10},
    table/col sep=comma
]
    % Загрузка данных из CSV, который создаётся при degradation_sweep
    \addplot[blue, thick, mark=*] table [
        x=Frequency_Hz, 
        y=THD_%
    ] {./data/thd_vs_freq.csv};
    \legend{THD}
\end{axis}
\end{tikzpicture}
\caption{Деградация THD от частоты. Данные получены с помощью автоматизированного свипирования.}
\label{fig:thd_vs_freq}
\end{figure}

Полученная зависимость подтверждает, что усилитель способен работать с высоким качеством (THD $< 0.06\%$) в полосе до 1\,МГц при выходной мощности $\sim 4.5$\,Вт. Дальнейшее расширение полосы возможно после оптимизации цепи коррекции.

% =============================================
% 6. РЕКОМЕНДАЦИИ
% =============================================
\section{Рекомендации по защите выхода}

\subsection{Защита выхода ADA4870}

При работе на индуктивную нагрузку возможны выбросы напряжения
за пределы питания. Рекомендуется установить защитные диоды
Шоттки непосредственно на выводе AOUT.

Схема включения (рис.~\ref{fig:diode_scheme}):
\begin{figure}[H]
\centering
\begin{circuitikz}[american, scale=1.1]
    \draw (0,0) node[op amp, noinv input up, label={[yshift=-1cm, xshift=-0.5cm]below:ADA4870}] (oa) {};
    \draw (oa.out) -- ++(0.5,0) coordinate (out_node);
    \draw (oa.+) node[ocirc] {};
    \draw (oa.-) node[ocirc] {};
    \draw (out_node) to[Do, l_=VD1, a^={MBRS360}, *-*] ++(0, 2)
          node[above] {$+18$~V (VDD)};
    \draw (out_node) to[Do, l_=VD2, a^={MBRS360}, *-*] ++(0,-2)
          node[below] {$-18$~V (VEE)};
    \draw (out_node) to[R, l=$R_{\text{Load}}$\ 22~Ом, *-o] ++(3,0)
          node[right] {$V_{\text{out}}$};
\end{circuitikz}
\caption{Схема включения защитных диодов на выходе ADA4870.}
\label{fig:diode_scheme}
\end{figure}

\begin{itemize}
  \item \textbf{VD1} (к положительной шине): анод к AOUT,
        катод к шине +18~В (VDD).
  \item \textbf{VD2} (к отрицательной шине): анод к шине $-18$~В
        (VEE), катод к AOUT.
\end{itemize}

\subsubsection{Выбор диода MBRS360}
MBRS360 \cite{MBRS360_ds} – диод Шоттки 3~А, 60~В в корпусе SMC.
Пиковый ток нагрузки 0{,}636~А значительно меньше предельного,
обратное напряжение 60~В с запасом перекрывает возможные выбросы.
Ёмкость диода (около 180~пФ) практически не влияет на сигнал
до 2~МГц.

\subsection{Компоненты для реальной платы}

\begin{itemize}
  \item \textbf{\(R_{Load}\)}: неиндуктивный резистор 22~Ом,
        мощностью не менее 10~Вт.
  \item \textbf{Защитные диоды}: MBRS360 (или аналоги с током \(\ge 1\)~А).
  \item \textbf{Байпасные конденсаторы}: 100~нФ + 10~мкФ
        от каждой шины питания к GND, вплотную к ADA4870.
  \item \textbf{\(C_f\)}: начинать с 2~пФ, при необходимости
        увеличить до 2{,}7--3{,}3~пФ по результатам AC-анализа.
  \item \textbf{Цепь TIA}: резисторы \(R_3=R_5=249\)~Ом,
        конденсаторы в цепи обратной связи ADA4807-2 не требуются
        (работа на низкоомную нагрузку).
\end{itemize}

% =============================================
% СПИСОК ЛИТЕРАТУРЫ
% =============================================
\begin{thebibliography}{99}

\bibitem{ADA4870_ds}
Analog Devices, Inc. \textit{ADA4870: High Speed, High Voltage,
1\,A Output Drive Amplifier. Data Sheet.} Rev.\,C, 2025.
URL: \url{https://www.analog.com/en/products/ada4870.html}.

\bibitem{ADA4807_ds}
Analog Devices, Inc. \textit{ADA4807-1/ADA4807-2/ADA4807-4:
Low Noise, Rail-to-Rail Amplifier. Data Sheet.} Rev.\,D, 2020.
URL: \url{https://www.analog.com}.

\bibitem{AD9705_ds}
Analog Devices, Inc. \textit{AD9704/AD9705/AD9706/AD9707:
8-/10-/12-/14-Bit, 175\,MSPS TxDAC Digital-to-Analog Converters.
Data Sheet.} Rev.\,G, 2023.
URL: \url{https://www.analog.com}.

\bibitem{MBRS360_ds}
ON Semiconductor. \textit{MBRS360T3G: Surface Mount Schottky Power
Rectifier. Data Sheet.} Rev.\,14, March 2020.
URL: \url{https://www.onsemi.com}.

\bibitem{LTspice}
Analog Devices, Inc. \textit{LTspice Simulation Software.}
Версия~26.0.1, 2026~г.

\end{thebibliography}

\end{document}